\section{Reconstrucción íntegra del registro e isometría de incidencia}
\label{rev4:ki:section}

La reconstrucción de la sección~\ref{sec:memory} permite conservar en
una misma composición la geometría positiva, la incidencia excepcional
y la dirección de deformación. La información inicial de esta
composición es el registro producido por el estado enriquecido
APP--TRIT--TPK: las dos hojas aritméticas conservan sus operaciones;
la orientación trítica determina el régimen; el transporte, la
actualización y la prolongación nonádica conservan las diferencias
entre posiciones y la componente uniforme. Las realizaciones lineales
que siguen actúan sobre ese registro ya determinado. Su función es
transportar información y explicitar sus invariantes.

La distinción fundamental se refiere al dominio conservado. El
registro \(K\) contiene doce coordenadas ordenadas; la tétrada
\(B_K\) conserva cuatro posiciones seleccionadas por la incidencia;
una dirección geométrica conserva una recta en un espacio de
representación. Cada operación tiene una imagen y una fibra propias.
La construcción siguiente conserva primero las doce coordenadas y
deriva después las reducciones que intervienen en la geometría. Así
queda disponible una reconstrucción exacta antes de efectuar una
extracción de soporte o una reducción dimensional.

\subsection{Forma de Gram del diseño dodecafásico}
\label{rev4:ki:gram-subsection}

Adoptamos la numeración de uno a doce y denotamos por
\(\Omega=\{1,\ldots,12\}\) el conjunto de posiciones. La
codificación Paley--Witt del estado enriquecido determina el conjunto
\(\mathcal H\) de sus \(132\) hexadas. Una hexada es un soporte
de seis posiciones de un vector de peso seis del código ternario
extendido. La construcción explícita del código y la demostración de
que cada quintupla pertenece a una única hexada se desarrollan en
el teorema~\ref{thm:lf-code} y la
proposición~\ref{prop:lf-witt}. Utilizamos aquí precisamente esa
incidencia, con las posiciones transportadas en el mismo orden que
las coordenadas del registro.

Sean \(V=\RR^{12}\) y \(W=\RR^{\mathcal H}\), equipados con
sus productos interiores euclídeos. La matriz de incidencia define
la aplicación
\begin{equation}
 \mathcal I:V\longrightarrow W,
 \qquad (\mathcal I k)_H=\sum_{p\in H}k_p,
 \qquad \mathcal I_{H,p}=\mathbf1_{\{p\in H\}}.
 \label{rev4:ki:incidence-map}
\end{equation}
El símbolo \(\mathbf1\in V\) designa el vector uniforme de doce
componentes. Definimos
\[
 \mu(k)=\frac1{12}\mathbf1^{\mathsf T}k,\qquad
 P_0=I_{12}-\frac1{12}\mathbf1\mathbf1^{\mathsf T},\qquad
 V_0=\mathbf1^\perp.
\]
El proyector \(P_0\) coincide con \(P_{11}\) de la
sección~\ref{sec:memory}; el subíndice cero señala en esta
presentación la condición de media nula. La descomposición ortogonal
es \(k=\mu(k)\mathbf1+P_0k\).

\begin{lemma}[Identidad de Gram y normalización incidencial]
\label{rev4:ki:gram}
La incidencia satisface
\begin{equation}
 \mathcal I^{\mathsf T}\mathcal I
       =36I_{12}+30\mathbf1\mathbf1^{\mathsf T}.
 \label{rev4:ki:gram-formula}
\end{equation}
Por consiguiente,
\[
 U_W=\frac16\mathcal I\big|_{V_0}:
 V_0\longrightarrow E_{\rm inc}:=\mathcal I(V_0)
\]
es una isometría sobre un subespacio de dimensión once.
\end{lemma}
\begin{proof}
Fijado un punto, las quintuplas que lo contienen son
\(\binom{11}{4}\). Cada hexada que contiene ese punto comprende
\(\binom54\) de ellas. La unicidad de la extensión de cada
quintupla da \(66\) hexadas por punto. Fijada una pareja de
puntos, el mismo argumento da
\(\binom{10}{3}/\binom43=30\) hexadas por pareja.
La entrada \((p,r)\) de
\(\mathcal I^{\mathsf T}\mathcal I\) cuenta exactamente las
hexadas que contienen simultáneamente \(p\) y \(r\). Sus
entradas diagonales son, por tanto, \(66\), y las restantes,
\(30\), lo que prueba \eqref{rev4:ki:gram-formula}.

Para \(z\in V_0\), el término uniforme se anula y
\[
 \left\|\frac16\mathcal I z\right\|^2
 =\frac1{36}z^{\mathsf T}
       \mathcal I^{\mathsf T}\mathcal I z
 =\|z\|^2.
\]
La inyectividad y la dimensión de la imagen se deducen de esta
igualdad. El factor seis queda así fijado por el Gram del diseño.
\end{proof}

El conteo de incidencia determina la escala de la isometría. Las
sumas sobre hexadas son observables lineales de las posiciones; sus
productos interiores contienen la multiplicidad con que cada
posición y cada pareja aparecen en el diseño. La eliminación del
modo uniforme separa esos dos niveles de conteo y deja el factor
cuadrático \(36\). Esta operación proporciona un enlace explícito
entre la representación de doce posiciones y su representación
sobre hexadas.

\subsection{Conservación de la media e inversión de las doce coordenadas}
\label{rev4:ki:full-subsection}

La componente uniforme se conserva mediante una coordenada escalar
adicional. Equipamos \(\RR\oplus E_{\rm inc}\) con el producto
interior
\[
 \langle(\mu,y),(\nu,z)\rangle_*=
 12\mu\nu+\langle y,z\rangle.
\]
El peso doce es la norma cuadrada de \(\mathbf1\), fijada por el
número de posiciones. Definimos
\begin{equation}
 F:V\longrightarrow\RR\oplus E_{\rm inc},\qquad
 F(k)=\left(\mu(k),\frac16\mathcal I P_0k\right).
 \label{rev4:ki:F}
\end{equation}

\begin{theorem}[Isomorfismo isométrico del registro completo]
\label{rev4:ki:full-isometry}
La aplicación \(F\) es un isomorfismo lineal isométrico. Su inversa
es
\begin{equation}
 F^{-1}(\mu,y)=\mu\mathbf1+
                  \frac16P_0\mathcal I^{\mathsf T}y,
 \qquad y\in E_{\rm inc}.
 \label{rev4:ki:F-inverse}
\end{equation}
En particular,
\begin{equation}
 12|\mu(k)|^2+
 \left\|\frac16\mathcal I P_0k\right\|^2=\|k\|^2.
 \label{rev4:ki:F-norm}
\end{equation}
\end{theorem}
\begin{proof}
La descomposición uniforme--centrada es ortogonal y el
lema~\ref{rev4:ki:gram} conserva la norma de su segundo sumando.
Esto demuestra \eqref{rev4:ki:F-norm} y la inyectividad.
Además,
\[
 \frac1{36}P_0\mathcal I^{\mathsf T}\mathcal I P_0=P_0.
\]
Al sustituir \(F(k)\) en \eqref{rev4:ki:F-inverse}, se obtiene
\(\mu(k)\mathbf1+P_0k=k\). Recíprocamente, si
\(y\in E_{\rm inc}\), existe un único \(z\in V_0\) con
\(y=\mathcal I z/6\). La fórmula propuesta devuelve
\(\mu\mathbf1+z\), cuya imagen es \((\mu,y)\). Así se
prueban la sobreyectividad sobre el codominio declarado y las dos
identidades de inversión.
\end{proof}

La restricción \(y\in E_{\rm inc}\) expresa una compatibilidad
lineal efectiva. El operador
\begin{equation}
 \Pi_{\rm inc}=\frac1{36}\mathcal I P_0
                         \mathcal I^{\mathsf T}:W\longrightarrow W
 \label{rev4:ki:inc-projector}
\end{equation}
es el proyector ortogonal sobre \(E_{\rm inc}\). En efecto, es
simétrico y su cuadrado se reduce a él mediante
\eqref{rev4:ki:gram-formula}; actúa como la identidad en la imagen
de \(\mathcal I P_0\). Por tanto, la compatibilidad se comprueba
mediante \(\Pi_{\rm inc}y=y\). Para un dato arbitrario de
\(W\), la reconstrucción seguida de incidencia devuelve
\(\Pi_{\rm inc}y\); el componente ortogonal
\((I_W-\Pi_{\rm inc})y\) cuantifica la información externa a
esta representación del registro.

\begin{proposition}[Covariancia de la reconstrucción]
\label{rev4:ki:equivariance}
Sea \(g\) un automorfismo del diseño y sean \(\rho(g)\) y
\(\sigma(g)\) sus acciones por permutación sobre posiciones y
hexadas. Entonces
\[
 F\rho(g)=(1\oplus\sigma(g))F.
\]
La media se conserva y la reconstrucción conmuta con el transporte
simultáneo de las posiciones y de sus hexadas.
\end{proposition}
\begin{proof}
La equivalencia \(p\in H\Longleftrightarrow g(p)\in g(H)\)
implica \(\mathcal I\rho(g)=\sigma(g)\mathcal I\).
Las permutaciones fijan \(\mathbf1\), conservan la media y
conmutan con \(P_0\). Sustituir estas identidades en
\eqref{rev4:ki:F} demuestra la afirmación. La covariancia de la
inversa se obtiene componiendo con \(F^{-1}\).
\end{proof}

La normalización isométrica y la inversión aritmética cumplen
funciones distintas. En la coordenatización Hadamard del registro,
\(U=\mathcal H_{12}K\) y
\(K=\mathcal H_{12}U/4\), con
\(\mathcal H_{12}^{\mathsf T}=\mathcal H_{12}\) y
\(\mathcal H_{12}^2=4I_{12}\).
El operador \(J_H=\mathcal H_{12}/2\) es ortogonal. Por ello
\(FJ_H\) conserva la norma del registro firmado normalizado
\(U/2\) y lo representa en el espacio incidencial ampliado.
La inversión integral mediante \(\mathcal H_{12}/4\) conserva,
además, la subred de congruencias correspondiente. En ambos casos se
transporta el registro generado, con el orden de fase fijado antes
de la transformación \cite{hmtX}.

\subsection{Composición exacta desde las dos memorias nonádicas}
\label{rev4:ki:memory-subsection}

Retomamos \(M\), \(q\), \(m_0\) y \(m_1\)
de \eqref{eq:mem-data}, junto con el operador de reconstrucción
\(L=(M^*M|_{V_0})^{-1}M^*\) de la sección~\ref{sec:memory}.
El dominio compatible es
\(\mathcal D=\RR\times\operatorname{im}M\), con
\(m=(m_0,m_1)\). El teorema~\ref{thm:memory-inversion}
establece la reconstrucción
\[
 \mathcal R(q,m)=\frac q{12}\mathbf1+Lm,
 \qquad LM=P_0.
\]
Como \(L\) toma valores en \(V_0\), la composición con la
incidencia tiene una expresión directa:
\begin{equation}
 \mathcal C=F\mathcal R,\qquad
 \mathcal C(q,m)=\left(\frac q{12},\frac16\mathcal I Lm\right).
 \label{rev4:ki:memory-incidence}
\end{equation}

\begin{corollary}[Equivalencia entre memoria compatible e incidencia íntegra]
\label{rev4:ki:composition}
La aplicación \(\mathcal C\) es un isomorfismo entre
\(\mathcal D\) y \(\RR\oplus E_{\rm inc}\), con inversa
\begin{equation}
 \mathcal C^{-1}(\mu,y)=
 \left(12\mu,\frac16MP_0\mathcal I^{\mathsf T}y\right).
 \label{rev4:ki:C-inverse}
\end{equation}
Su forma cuadrática conservada es
\[
 \|\mathcal C(q,m)\|_*^2=
 \frac{|q|^2}{12}+\|Lm\|^2=\|\mathcal R(q,m)\|^2.
\]
\end{corollary}
\begin{proof}
La fórmula \eqref{rev4:ki:memory-incidence} resulta de
\(P_0\mathbf1=0\) y \(P_0L=L\). La inversa se obtiene
aplicando \(k\mapsto(\mathbf1^{\mathsf T}k,Mk)\) a
\eqref{rev4:ki:F-inverse}; se usa \(M\mathbf1=0\).
Las identidades de composición proceden de \(LM=P_0\) y de
\(MLm=m\) para \(m\in\operatorname{im}M\). La igualdad de
normas es \eqref{rev4:ki:F-norm} aplicada a \(\mathcal R(q,m)\).
\end{proof}

\begin{figure}[htbp]
\centering
\begin{tikzpicture}[>=Latex, every node/.style={font=\small},
  box/.style={draw,rounded corners=2pt,align=center,
              inner sep=6pt,minimum height=12mm}]
 \node[box] (memory) at (0,0) {Carga y memorias\\$(q,m_0,m_1)$};
 \node[box] (register) at (4.1,0) {Registro íntegro\\$K$};
 \node[box] (incidence) at (8.2,0) {Media e incidencia\\$(\mu,y)$};
 \node[box] (direction) at (4.1,-1.75)
    {Dirección geométrica\\$\ell_K=\RR P_3K$};
 \draw[<->] (memory) -- node[above] {$\mathcal R$} (register);
 \draw[<->] (register) -- node[above] {$F$} (incidence);
 \draw[->] (register) -- node[right] {proyección y paso a la recta}
    (direction);
\end{tikzpicture}
\caption{La carga y las dos memorias compatibles reconstruyen las doce
coordenadas. La media y el componente incidencial proporcionan otra
representación reversible del mismo registro. La dirección geométrica
es una reducción posterior con fibras de dimensión positiva.}
\label{rev4:ki:figure}
\end{figure}

Esta composición determina también la partición positiva de
\ref{sec:memory}. Si \(q>0\), sus longitudes se recuperan mediante
\begin{equation}
 d_j=\frac1{12}+\frac{1}{6q}
             \bigl(P_0\mathcal I^{\mathsf T}y\bigr)_j,
 \qquad y=\frac16\mathcal I P_0K.
 \label{rev4:ki:positive-lengths}
\end{equation}
Las doce condiciones \(d_j>0\) son las ecuaciones de la coordenatización
positiva expresadas en coordenadas de incidencia. Con
\(t_j=\sum_{r\leq j}d_r\), los menores ordenados
\(t_j-t_i=\sum_{i<r\leq j}d_r\) se reconstruyen desde
\((\mu,y)\) o desde \((q,m_0,m_1)\). De este modo, la
positividad de los menores y la conservación isométrica de la
incidencia quedan relacionadas por aplicaciones explícitas. La
primera es una desigualdad en la coordenatización ordenada; la segunda es una
identidad de productos interiores.

\subsection{Recuperación de la polarización y alcance de las reducciones}
\label{rev4:ki:direction-subsection}

La realización de las doce posiciones sobre las clases de
\(A_5/C_5\) determina el proyector tridimensional
\[
 P_3=\frac3{60}\sum_{a\in A_5}
               \chi_3(a^{-1})\varrho(a).
\]
Aquí \(\varrho\) es la representación por permutaciones y
\(\chi_3\) el carácter tridimensional elegido en la coordenatización
marcada. Sus representantes, valores de carácter y matriz efectiva
se fijan en \eqref{eq:leech-projectors}; el
lema~\ref{lem:leech-u} demuestra
\(P_3=P_3^{\mathsf T}=P_3^2\),
\(\operatorname{rank}P_3=3\) y \(P_3\mathbf1=0\).
En esa coordenatización, la dirección del registro se obtiene por
\begin{equation}
 u_K=P_3P_0K=P_3Lm
     =\frac16P_3P_0\mathcal I^{\mathsf T}y.
 \label{rev4:ki:direction-reconstruction}
\end{equation}
Las igualdades siguen de las dos inversas ya demostradas. Por tanto,
la misma dirección se recupera desde las memorias o desde el
componente incidencial; la media se conserva para reconstruir
\(K\), aunque el proyector tridimensional la anule.

Para el registro de \eqref{eq:leech-K}, la evaluación exacta
\eqref{eq:leech-u-norm} da
\[
 \|u_K\|^2=
       \frac{6638585+2275584\sqrt5}{20}>0.
\]
La recta \(\ell_K=\RR u_K\) y el complemento transversal
\(V_{10}(K)=V_0\cap u_K^\perp\) quedan así determinados.
Su proyector es
\begin{equation}
 P_{10}(K)=P_0-
       \frac{u_Ku_K^{\mathsf T}}{\|u_K\|^2}.
 \label{rev4:ki:rank-ten}
\end{equation}
En efecto, el segundo sumando es un proyector ortogonal de rango uno
contenido en \(V_0\); su producto con \(P_0\), en cualquiera
de los dos órdenes, coincide consigo mismo. Al elevar
\eqref{rev4:ki:rank-ten} al cuadrado se recupera el mismo
operador, y su traza vale \(11-1=10\). Se obtiene
\(V_0=\ell_K\oplus^\perp V_{10}(K)\).
Esta reducción es la que interviene después en la polarización de
los doce planos reticulares.

La incidencia transporta asimismo estos operadores. Sobre
\(E_{\rm inc}\), el operador
\(\widehat P_3=U_WP_3U_W^*\) satisface
\(\widehat P_3U_W=U_WP_3\); la identidad procede de
\(U_W^*U_W=I_{V_0}\). Transportar de igual modo
\(P_{10}(K)\) conserva su rango y su producto interior.
La correspondencia entre las dos representaciones se establece,
por tanto, sobre vectores y operadores, además de sus dimensiones.

La necesidad de conservar los distintos componentes admite
comprobaciones precisas. La observación centrada de \(k\) coincide
con la de \(k+c\mathbf1\) para todo \(c\), mientras sus
cargas difieren en \(12c\). En particular, \(k=0\) y
\(k=\mathbf1\) tienen el mismo componente incidencial centrado
y medias distintas. La coordenada \(\mu\) es exactamente la
información que separa esos dos registros. Del mismo modo,
\(u_{ak+b\mathbf1}=a u_k\) para \(a\ne0\); la recta
resultante coincide con \(\ell_k\), aunque cambien la amplitud
y la media. Estas identidades describen las fibras de las
aplicaciones sobre el espacio ambiente declarado.

La extracción \(k\mapsto\{j:k_j\geq729\}\) posee otras
fibras. Para el registro concreto \(K\), tanto \(K\) como
\(K+e_1\) tienen el soporte \(\{4,6,9,10\}\), donde
\(e_1\) es el primer vector de la base canónica. Son dos vectores
ambiente diferentes con la misma tétrada. Esta comprobación
caracteriza la pérdida de información del mapa de soporte; la
selección terminal del registro permanece fijada por su
construcción previa. La bandera
\(B_K\subset H^-_\alpha\) conserva además la hexada determinada
por los signos anteriores al acarreo, cuya coincidencia con la
selección incidencial se demuestra en la
proposición~\ref{prop:lf-flag}.

El resultado común de estas composiciones es una continuidad
matemática explícita. La memoria compatible reconstruye el registro;
su descomposición uniforme--centrada determina una representación
incidencial isométrica; la reconstrucción de sus longitudes determina
la coordenatización positiva; y la proyección tridimensional determina la
dirección de deformación. La tétrada y la bandera conservan las
incidencias que utilizan las construcciones excepcionales, mientras
\(F\) conserva la totalidad del registro dodecafásico. La historia
enriquecida mantiene, a su vez, la orientación, el recorrido y la
prolongación que sitúan ese registro dentro de la estructura
discreta conjunta del continuo.
